\section{هفته‌ی 3: برازش مدل با کمترین مربعات: سطح خطا و بهینه‌سازی}

\begin{quote}
\textbf{فصل 2: روش‌های برآورد پارامتر} · درس 1 ساعت · \lr{lab} 2 ساعت
\textbf{داده‌ی اصلی این هفته:} آزمایش فراموشی شبیه‌سازی‌شده از هفته‌ی 2 (مدل
توانی، \(A^\ast = 0.9\)، \(d^\ast = 0.5\)، \(n = 50\) آیتم در 8 فاصله).
آن را درست برازش می‌کنیم و با برازش‌های خطی‌شده‌ی هفته‌ی قبل مقایسه می‌کنیم.
\end{quote}

\begin{center}\rule{0.5\linewidth}{0.5pt}\end{center}

\subsection{واژه‌های کلیدی}

واژه‌های مهم این درس با معادل فارسی‌شان (حتی اگر در درس‌های قبل آمده
باشند).

{\def\LTcaptype{none} % do not increment counter
\begin{longtable}[]{@{}
  >{\raggedright\arraybackslash}p{(\linewidth - 4\tabcolsep) * \real{0.3333}}
  >{\raggedright\arraybackslash}p{(\linewidth - 4\tabcolsep) * \real{0.3333}}
  >{\raggedright\arraybackslash}p{(\linewidth - 4\tabcolsep) * \real{0.3333}}@{}}
\toprule\noalign{}
\begin{minipage}[b]{\linewidth}\raggedright
\lr{English}
\end{minipage} & \begin{minipage}[b]{\linewidth}\raggedright
فارسی
\end{minipage} & \begin{minipage}[b]{\linewidth}\raggedright
توضیح کوتاه
\end{minipage} \\
\midrule\noalign{}
\endhead
\bottomrule\noalign{}
\endlastfoot
\lr{fit} / \lr{fitting} & برازش & تنظیم پارامترها تا مدل با داده جور
شود \\
\lr{least squares} & کمترین مربعات & روش برازش با کمینه کردن مجموع مربع
خطاها \\
\lr{error surface} & سطح خطا & خطا به‌عنوان تابعی از پارامترها \\
\lr{optimizer} & بهینه‌ساز & الگوریتمی که کمینه‌ی تابع را پیدا می‌کند \\
\lr{parameter} & پارامتر & کمیت قابل تنظیم مدل که بین افراد یا شرایط فرق
می‌کند \\
\lr{estimate} / \lr{estimator} & برآورد / برآوردگر & مقدار به‌دست‌آمده
برای پارامتر / روش به دست آوردن آن \\
\lr{forgetting curve} & منحنی فراموشی & کاهش یادآوری با گذشت زمان \\
\lr{exponential} / \lr{power function} & تابع نمایی / توانی & دو شکل
رایج برای منحنی فراموشی و یادگیری \\
\lr{SSE} / \lr{RMSD} & مجموع مربعات خطا / ریشه‌ی میانگین مربع انحراف &
معیارهای فاصله‌ی داده و پیش‌بینی \\
\lr{noise} & نویز & تغییرات تصادفی در شواهد یا در داده \\
\lr{parameter trade-off} & بده‌بستان پارامترها & تغییر یک پارامتر که با
پارامتر دیگر جبران می‌شود \\
\lr{grid approximation} & تقریب شبکه‌ای \lr{(grid)} & محاسبه‌ی پسین روی
شبکه‌ای از مقادیر \\
\lr{Nelder\lr{--}Mead} & روش \lr{Nelder\lr{--}Mead} & بهینه‌ساز بدون نیاز به
مشتق \\
\lr{Bayes factor (BF)} & عامل بیز \lr{(Bayes factor)} & نسبت
درست‌نمایی‌های حاشیه‌ای دو مدل \\
\lr{local minimum} & کمینه‌ی محلی & نقطه‌ای پایین‌تر از اطراف ولی نه بهترین
نقطه \\
\lr{bounds} / \lr{constraints} & کران‌ها / قیدها & محدودیت‌های مجاز برای
مقدار پارامترها \\
\lr{reparameterization} & بازپارامتری‌سازی & نوشتن پارامتر در مقیاس دیگر،
مثلاً \ensuremath{\lambda} = \lr{e\^{}\ensuremath{\eta}} \\
\lr{multi-start} & شروع چندگانه \lr{(multi-start)} & اجرای بهینه‌ساز از
چند نقطه‌ی شروع \\
\lr{maximum likelihood estimate (MLE)} & برآورد بیشینه‌ی درست‌نمایی &
پارامترهایی که داده را محتمل‌ترین می‌کنند \\
\lr{likelihood} & درست‌نمایی & احتمال داده به‌عنوان تابعی از پارامترها \\
\lr{gradient} & گرادیان & جهت و شیب تغییر تابع \\
\lr{convergence} & همگرایی & رسیدن زنجیره‌ها به توزیع هدف \\
\lr{large language model (LLM)} & مدل بزرگ زبانی & شبکه‌ی عصبی‌ای که کلمه‌ی
بعدی متن را پیش‌بینی می‌کند \\
\lr{scaling law} & قانون مقیاس & رابطه‌ی توانی بین اندازه‌ی مدل و خطا \\
\end{longtable}
}

\begin{center}\rule{0.5\linewidth}{0.5pt}\end{center}

\subsection{اهداف یادگیری}

در پایان این هفته، دانشجویان می‌توانند:

\begin{enumerate}
\def\labelenumi{\arabic{enumi}.}
\tightlist
\item
  \lr{SSE}، \lr{MSE} و \lr{RMSD} را به‌عنوان توابع ناهمخوانی
  \lr{(discrepancy functions)} \textbf{تعریف کنند} و \textbf{توضیح دهند}
  که کمینه کردن آن‌ها چه فرضی درباره‌ی نویز دارد.
\item
  سطح خطای یک مدل دوپارامتری را \textbf{رسم کنند} و شکل آن (کمینه، دره،
  فلات) را بر حسب \lr{trade-off} پارامترها \textbf{تفسیر کنند}.
\item
  برازش کمترین مربعات را با جستجوی شبکه‌ای و با
  \lr{\nolinkurl{scipy.optimize.minimize}} \lr{(Nelder\lr{--}Mead} و
  \lr{L-BFGS-B)} \textbf{پیاده‌سازی کنند} و نقاط قوت و حالت‌های شکست
  الگوریتم‌ها را \textbf{مقایسه کنند}.
\item
  شکست‌های رایج برازش را \textbf{تشخیص دهند} و \textbf{رفع کنند}: مقدار
  شروع بد، فلات‌ها، کمینه‌های محلی، و پارامترهای خارج از دامنه (با کران،
  \lr{reparameterization} و \lr{multi-start).}
\item
  یک برازش کمترین مربعات را با \lr{RMSD} و باقی‌مانده‌ها \textbf{ارزیابی
  کنند} و پارامترهای برازش‌شده را \textbf{تفسیر کنند}.
\end{enumerate}

\begin{center}\rule{0.5\linewidth}{0.5pt}\end{center}

\subsection{مقدمه}

هفته‌ی قبل منحنی‌های فراموشی را با کشیدن یک خط راست از میان نسبت‌های
لگاریتمی‌شده برازش کردیم. مدل توانی مدل واقعی بود، ولی برآوردهایش آشکارا
خطا داشتند (\(\hat A = 0.80\)، \(\hat d = 0.37\) به‌جای \(0.9\) و
\(0.5\)). این روش هر وقت یک شمارش صفر باشد از کار می‌افتد، و خطاها را روی
مقیاس غلط (لگاریتمی) کمینه می‌کند.

راه‌حل از نظر مفهومی ساده است: \textbf{پارامترهایی را انتخاب کن که
پیش‌بینی‌های مدل را، روی مقیاس خود داده، تا جای ممکن به داده نزدیک کنند.}
«تا جای ممکن نزدیک» به یک عدد نیاز دارد: یک \emph{تابع ناهمخوانی}.
«پارامترها را انتخاب کن» به یک رویه‌ی جستجو نیاز دارد: یک \emph{الگوریتم
بهینه‌سازی}. این هفته هر دو را پوشش می‌دهد.

دلیل عمیق‌تری هم هست که این برای علوم شناختی مهم است. پارامترهای مدل
\emph{اندازه‌گیری‌های} ما از سازه‌های پنهان‌اند (نرخ زوال، احتیاط، حساسیت).
رویه‌ی برازش ضعیف اندازه‌گیری‌های ضعیف می‌دهد، و هر نتیجه‌ای که از آن‌ها گرفته
شود (تفاوت گروه‌ها، همبستگی با فعالیت مغز، مقایسه‌ی مدل‌ها) آن خطا را به
ارث می‌برد. بدتر اینکه، همان‌طور که خواهیم دید، یک بهینه‌ساز می‌تواند
«موفقیت» گزارش کند در حالی که جواب بی‌معنی برمی‌گرداند.

\paragraph{این موضوع کجا به کار می‌رود؟}

\begin{itemize}
\tightlist
\item
  \textbf{هر مطالعه‌ی مبتنی بر مدل در روان‌شناسی} یک مرحله‌ی برازش دارد.
  منحنی‌های یادگیری، منحنی‌های فراموشی، تابع‌های روان‌سنجی \lr{(psychometric
  function:} اینکه تشخیص با قوی‌تر شدن محرک چطور بهتر می‌شود) و منحنی‌های
  دوز\lr{--}پاسخ در روان‌داروشناسی همه با کمینه کردن یک معیار خطا برازش
  می‌شوند.
\item
  \textbf{پژوهش‌های بالینی و رشدی} پارامترهای برازش‌شده را بین گروه‌ها
  مقایسه می‌کنند (مثلاً نرخ فراموشی بیماران در برابر گروه کنترل). اگر
  برازش بد باشد یا روی یک کران گیر کرده باشد، مقایسه‌ی گروه‌ها غلط است.
\item
  \textbf{علوم اعصاب} منحنی‌های تنظیم \lr{(tuning curve)} نورون‌ها و تابع
  پاسخ همودینامیک در \lr{fMRI} را با همین بهینه‌سازها برازش می‌کند.
\item
  \textbf{یادگیری ماشین} مدل‌ها را با کمینه کردن یک تابع زیان \lr{(loss)}
  با روش‌های گرادیانی آموزش می‌دهد؛ واژه‌ها (سطح زیان، نرخ یادگیری، کمینه‌ی
  محلی) همان واژه‌های این هفته‌اند.
\end{itemize}

\begin{center}\rule{0.5\linewidth}{0.5pt}\end{center}

\subsection{جزوه‌ی درس}

\paragraph{به زبان ساده}

مدل چند پیچ تنظیم (پارامتر) دارد. برازش یعنی چرخاندن پیچ‌ها تا منحنی مدل
تا جای ممکن از نزدیک نقطه‌های داده بگذرد. برای اینکه کامپیوتر این کار را
بکند دو چیز لازم داریم. اول، یک عدد که بگوید برازش فعلی چقدر بد است:
\textbf{مجموع مربعات خطا \lr{(SSE)}}، یعنی جمع مربع فاصله‌ی عمودی نقطه‌های
داده از منحنی. دوم، روشی برای چرخاندن پیچ‌ها که این عدد را کم کند: یک
\textbf{بهینه‌ساز \lr{(optimizer)}}. \lr{SSE} را مثل یک منظره روی همه‌ی
تنظیم‌های ممکن پیچ‌ها تصور کنید، با تپه‌ها (برازش بد) و دره‌ها (برازش خوب)؛
بهینه‌ساز یک کوهنورد است که دنبال پایین‌ترین نقطه می‌گردد. مشکل وقتی پیش
می‌آید که منظره یک دره‌ی دراز و صاف داشته باشد (دو پیچ می‌توانند اثر هم را
جبران کنند)، یک فلات صاف (کوهنورد نمی‌فهمد کدام طرف سرازیری است)، یا چند
دره‌ی جدا (کوهنورد ممکن است در دره‌ی اشتباه بایستد). بیشتر این هفته
درباره‌ی تشخیص این موقعیت‌ها و محافظت از خودمان در برابر آن‌هاست.

\textbf{زمان‌بندی پیشنهادی (60 دقیقه):} توابع ناهمخوانی 10 · سطح خطا 10 ·
الگوریتم‌های بهینه‌سازی 20 · مشکلات عملی و چک‌لیست برازش 15 · پیش‌نمایش 5.

\subsubsection{1. توابع ناهمخوانی}

برای مشاهدات \(y_i\) و پیش‌بینی‌های مدل
\(\hat y_i = f(x_i; \boldsymbol\theta)\)، \(i = 1, \dots, N\):

\[
\mathrm{SSE}(\boldsymbol\theta) = \sum_{i=1}^{N} \big(y_i - f(x_i;\boldsymbol\theta)\big)^2, \qquad
\mathrm{MSE} = \frac{\mathrm{SSE}}{N}, \qquad
\mathrm{RMSD} = \sqrt{\mathrm{MSE}}.
\]

\begin{itemize}
\tightlist
\item
  هر سه \textbf{کمینه‌کننده‌ی یکسانی} دارند: تبدیل‌های یکنوای همدیگرند.
  \lr{SSE} را \emph{کمینه} و \lr{RMSD} را \emph{گزارش} می‌کنیم، چون
  \lr{RMSD} بر حسب واحد داده است (اینجا: نسبت یادآوری)، پس \lr{«RMSD} =
  0.057» یعنی پیش‌بینی‌ها به‌طور میانگین حدود 0.06 خطا دارند.
\item
  \textbf{برآورد کمترین مربعات} این است: \[
  \hat{\boldsymbol\theta}_{\text{\lr{LS}}} = \arg\min_{\boldsymbol\theta} \mathrm{SSE}(\boldsymbol\theta).
  \]
\end{itemize}

\textbf{چرا مربع؟} مربع کردن همه‌ی خطاها را مثبت می‌کند، خطاهای بزرگ را
به‌شدت جریمه می‌کند، و تابعی هموار (مشتق‌پذیر) می‌دهد که بهینه‌سازها راحت با
آن کار می‌کنند. در سطحی عمیق‌تر، کمینه کردن \lr{SSE} دقیقاً همان برآورد
بیشینه درست‌نمایی است، \textbf{اگر} نویز گاوسی و با واریانس یکسان برای
همه‌ی مشاهدات باشد (هفته‌ی 4). این «اگر» فرض پنهان کمترین مربعات است.

\textbf{وقتی فرض برقرار نیست.} نسبت‌های یادآوری دوجمله‌ای‌اند: واریانسشان
\(f(1-f)/n\) به \(f\) بستگی دارد. نقاط نزدیک 0.5 پرنویزتر از نقاط نزدیک
0 یا 1 هستند، ولی \lr{SSE} ساده با همه‌شان مثل هم رفتار می‌کند.
\textbf{کمترین مربعات وزنی} \lr{(weighted least squares)} این را با وزن
دادن به هر مربع خطا بر اساس معکوس واریانسش درست می‌کند:

\[
\mathrm{WSSE}(\boldsymbol\theta) = \sum_{i} \frac{\big(y_i - \hat y_i\big)^2}{\mathrm{Var}(y_i)}.
\]

این را در تمرین 3 امتحان می‌کنید. هفته‌ی بعد، بیشینه درست‌نمایی این مسئله
را طبیعی‌تر حل می‌کند.

\begin{quote}
\textbf{نکته:} بعضی کتاب‌ها به \lr{RMSD} می‌گویند \lr{«RMSE».} هر دو یک
کمیت‌اند.
\end{quote}

\paragraph{خواندن ریاضی}

\begin{itemize}
\tightlist
\item
  \(\sum_{i=1}^{N}\) یعنی «روی همه‌ی \(N\) نقطه‌ی داده جمع بزن».
  \(y_i - f(x_i;\boldsymbol\theta)\) \textbf{باقی‌مانده \lr{(residual)}}
  نقطه‌ی \(i\) است: مقدار مشاهده‌شده منهای مقدار پیش‌بینی‌شده. مربع کردن و
  جمع زدن برای هر تنظیم پارامتر \(\boldsymbol\theta\) یک عدد می‌دهد.
\item
  \(\arg\min_{\boldsymbol\theta} \mathrm{SSE}(\boldsymbol\theta)\)
  این‌طور خوانده می‌شود: «مقداری از \(\boldsymbol\theta\) که در آن
  \lr{SSE} کمترین است». به تفاوت دقت کنید: \(\min\) کوچک‌ترین \emph{مقدار
  \lr{SSE}} است و \(\arg\min\) \emph{مقدار پارامتری} که به آن می‌رسد. ما
  \(\arg\min\) را می‌خواهیم.
\item
  \textbf{مثال کوچک.} داده‌ی \(y = (0.62, 0.41)\) و پیش‌بینی
  \((0.66, 0.44)\): باقی‌مانده‌ها \((-0.04, -0.03)\)، \lr{SSE}
  \(= 0.0016 + 0.0009 = 0.0025\)، \lr{MSE} \(= 0.00125\) و \lr{RMSD}
  \(= 0.035\): منحنی به‌طور میانگین حدود 3.5 درصد از نقطه‌ها فاصله دارد.
\item
  در \lr{WSSE} هر مربع باقی‌مانده بر واریانس همان نقطه تقسیم می‌شود، پس
  نقطه‌ی پرنویز وزن کمتری دارد. برای داده‌ی دوجمله‌ای واریانس برابر
  \(f(1-f)/n\) است.
\end{itemize}

\subsubsection{\texorpdfstring{2. سطح خطا \lr{(error
surface)}}{2. سطح خطا (error surface)}}

برای مدلی با \(k\) پارامتر، \(\mathrm{SSE}(\boldsymbol\theta)\) یک سطح
روی فضای پارامتر \(k\) بعدی است. با \(k = 2\) می‌توانیم آن را به شکل
نقشه‌ی تراز \lr{(contour map)} رسم کنیم: مسئله‌ی برازش می‌شود «پایین‌ترین
نقطه‌ی این منظره را پیدا کن.»

ویژگی‌هایی که باید بشناسید:

{\def\LTcaptype{none} % do not increment counter
\begin{longtable}[]{@{}
  >{\raggedright\arraybackslash}p{(\linewidth - 4\tabcolsep) * \real{0.3333}}
  >{\raggedright\arraybackslash}p{(\linewidth - 4\tabcolsep) * \real{0.3333}}
  >{\raggedright\arraybackslash}p{(\linewidth - 4\tabcolsep) * \real{0.3333}}@{}}
\toprule\noalign{}
\begin{minipage}[b]{\linewidth}\raggedright
ویژگی
\end{minipage} & \begin{minipage}[b]{\linewidth}\raggedright
ظاهرش
\end{minipage} & \begin{minipage}[b]{\linewidth}\raggedright
معنایش
\end{minipage} \\
\midrule\noalign{}
\endhead
\bottomrule\noalign{}
\endlastfoot
\textbf{یک کاسه‌ی گرد} & خطوط تراز هم‌مرکز و تقریباً دایره‌ای & پارامترها
به‌خوبی تعیین می‌شوند و تقریباً مستقل‌اند؛ هر الگوریتم معقولی کمینه را پیدا
می‌کند \\
\textbf{دره‌ی بلند و مورب} & خطوط تراز کشیده و کج &
\textbf{\lr{trade-off} پارامترها}: افزایش یک پارامتر را می‌شود با تغییر
دیگری جبران کرد. برآوردها در امتداد دره همبسته و نامطمئن خواهند بود
(هفته‌ی 5، هفته‌ی 7) \\
\textbf{فلات \lr{(plateau)}} & خطوط تراز خیلی دور از هم & پیش‌بینی‌های مدل
با پارامترها تقریباً تغییر نمی‌کنند (مثلاً نرخ زوال آن‌قدر بزرگ که منحنی
پیش‌بینی‌شده همه‌جا 0 است). روش‌های گرادیانی گیر می‌کنند \\
\textbf{چند حوضه‌ی جدا} & بیش از یک مجموعه‌ی بسته از خطوط تراز &
\textbf{کمینه‌های محلی}: جواب الگوریتم به نقطه‌ی شروع بستگی دارد \\
\end{longtable}
}

در فراموشی نمایی، \(A\) و \(\lambda\) با هم \lr{trade-off} دارند: سطح
شروع بالاتر با زوال سریع‌تر می‌تواند تقریباً همان مقادیر میانی را پیش‌بینی
کند که سطح پایین‌تر با زوال کندتر. به همین دلیل دره مورب است (در \lr{lab}
می‌بینید).

\subsubsection{3. الگوریتم‌های بهینه‌سازی}

همه‌ی الگوریتم‌های زیر از مقادیر اولیه‌ی \(\boldsymbol\theta^{(0)}\) شروع
می‌کنند و دنباله‌ای
\(\boldsymbol\theta^{(1)},\allowbreak \boldsymbol\theta^{(2)},\allowbreak \dots\) می‌سازند که
(امیدواریم) به کمینه پایین برود.

\paragraph{\texorpdfstring{3.1 جستجوی شبکه‌ای \lr{(grid
search)}}{3.1 جستجوی شبکه‌ای (grid search)}}

\lr{SSE} را در همه‌ی نقاط یک شبکه‌ی منظم حساب کنید و بهترین را نگه دارید.

\begin{itemize}
\tightlist
\item
  \textbf{مزایا:} فوق‌العاده ساده؛ گول کمینه‌های محلی \emph{روی شبکه} را
  نمی‌خورد؛ تصویری از کل سطح می‌دهد.
\item
  \textbf{معایب:} دقت محدود به فاصله‌ی نقاط شبکه است؛ هزینه با تعداد
  پارامترها \textbf{به‌طور نمایی} بالا می‌رود. با 100 مقدار برای هر
  پارامتر: \(k = 2\) به \(10^4\) ارزیابی و \(k = 5\) به \(10^{10}\)
  ارزیابی نیاز دارد. این همان \textbf{نفرین ابعاد} \lr{(curse of
  dimensionality)} است.
\item
  \textbf{کاربرد:} رسم سطح‌های 1 یا 2 پارامتری، و تولید مقادیر شروع خوب.
\end{itemize}

\paragraph{\texorpdfstring{3.2 سیمپلکس \lr{Nelder\lr{--}Mead} (بدون
مشتق)}{3.2 سیمپلکس Nelder\lr{--}Mead (بدون مشتق)}}

یک \textbf{سیمپلکس} مجموعه‌ای از \(k+1\) نقطه در فضای \(k\) بعدی است
(برای \(k = 2\) یک مثلث). هر تکرار سعی می‌کند \textbf{بدترین} رأس
\(\boldsymbol\theta_w\) را با حرکت دادنش از میان مرکز ثقل
\(\bar{\boldsymbol\theta}\) بقیه‌ی رأس‌ها جایگزین کند:

{\def\LTcaptype{none} % do not increment counter
\begin{longtable}[]{@{}
  >{\raggedright\arraybackslash}p{(\linewidth - 4\tabcolsep) * \real{0.3333}}
  >{\raggedright\arraybackslash}p{(\linewidth - 4\tabcolsep) * \real{0.3333}}
  >{\raggedright\arraybackslash}p{(\linewidth - 4\tabcolsep) * \real{0.3333}}@{}}
\toprule\noalign{}
\begin{minipage}[b]{\linewidth}\raggedright
حرکت
\end{minipage} & \begin{minipage}[b]{\linewidth}\raggedright
فرمول (ضرایب استاندارد)
\end{minipage} & \begin{minipage}[b]{\linewidth}\raggedright
چه وقت
\end{minipage} \\
\midrule\noalign{}
\endhead
\bottomrule\noalign{}
\endlastfoot
بازتاب \lr{(reflect)} &
\(\boldsymbol\theta_r = \bar{\boldsymbol\theta} + (\bar{\boldsymbol\theta} - \boldsymbol\theta_w)\)
& همیشه اول امتحان می‌شود \\
گسترش \lr{(expand)} &
\(\boldsymbol\theta_e = \bar{\boldsymbol\theta} + 2(\bar{\boldsymbol\theta} - \boldsymbol\theta_w)\)
& نقطه‌ی بازتاب‌شده از بهترین فعلی بهتر است؛ جلوتر برو \\
انقباض \lr{(contraction)} &
\(\boldsymbol\theta_c = \bar{\boldsymbol\theta} + \tfrac{1}{2}(\boldsymbol\theta_w - \bar{\boldsymbol\theta})\)
(داخلی)، یا همین گام در سمت بازتاب (خارجی) & نقطه‌ی بازتاب‌شده از دومین
بدترین رأس بهتر نیست \\
جمع‌شدن \lr{(shrink)} & همه‌ی رأس‌ها را تا نیمه به سمت بهترین ببر & انقباض
هم شکست خورد \\
\end{longtable}
}

نقطه‌ی بازتاب‌شده \textbf{پذیرفته می‌شود} وقتی از دومین بدترین رأس بهتر
باشد ولی از بهترین بهتر نباشد. سیمپلکس سرازیر «غلت می‌خورد»، در امتداد
دره‌ها کشیده می‌شود و دور کمینه جمع می‌شود.

\begin{itemize}
\tightlist
\item
  \textbf{مزایا:} فقط به مقادیر \lr{SSE} نیاز دارد (بدون مشتق)؛ در برابر
  توابع هدف پرنویز یا ناهموار مقاوم است؛ ابزار پیش‌فرض و پرکاربرد در
  مدل‌سازی شناختی.
\item
  \textbf{معایب:} در ابعاد بالا ممکن است کند باشد؛ در دره‌های بلند و
  باریک ممکن است گیر کند؛ هیچ تضمینی برای پیدا کردن کمینه‌ی سراسری ندارد.
\end{itemize}

\begin{quote}
روش‌های جستجوی مستقیم مثل \lr{Nelder\lr{--}Mead} در کتاب قدوسیان درباره‌ی
روش‌های مستقیم در بهینه‌سازی غیرخطی به‌تفصیل بررسی شده‌اند (مطالعه‌ی بیشتر را
ببینید).
\end{quote}

\paragraph{3.3 روش‌های گرادیانی}

\textbf{گرادیان} \(\nabla \mathrm{SSE}(\boldsymbol\theta)\) به سمت
سربالایی اشاره می‌کند، پس ما خلافش گام برمی‌داریم:

\[
\boldsymbol\theta^{(m+1)} = \boldsymbol\theta^{(m)} - \eta\,\nabla \mathrm{SSE}\big(\boldsymbol\theta^{(m)}\big),
\qquad
\frac{\partial\,\mathrm{SSE}}{\partial \theta} = -2 \sum_{i} \big(y_i - \hat y_i\big)\,\frac{\partial \hat y_i}{\partial \theta},
\]

با اندازه‌ی گام (نرخ یادگیری) \(\eta\). برای مدل نمایی،
\(\partial \hat y/\partial A = e^{-\lambda t}\) و
\(\partial \hat y/\partial \lambda = -A\,t\,e^{-\lambda t}\).

در عمل از روش‌های \textbf{شبه‌نیوتنی} \lr{(quasi-Newton)} مثل
\textbf{\lr{BFGS}} استفاده می‌کنیم. این روش‌ها \emph{انحنای} سطح را هم
تخمین می‌زنند و اندازه‌ی گام را خودکار انتخاب می‌کنند، پس نزدیک کمینه خیلی
سریع‌تر همگرا می‌شوند. \textbf{\lr{L-BFGS-B}} نسخه‌ای است که \textbf{کران}
پارامترها را می‌پذیرد. اگر گرادیان را ندهید، \lr{scipy} آن را عددی با
تفاضل‌های متناهی تقریب می‌زند.

\begin{itemize}
\tightlist
\item
  \textbf{مزایا:} روی سطوح هموار سریع و دقیق‌اند؛ برای تعداد پارامتر زیاد
  خوب مقیاس‌پذیرند (شبکه‌های عصبی همین‌طور آموزش می‌بینند).
\item
  \textbf{معایب:} به تابع هدف هموار نیاز دارند؛ هر جا گرادیان (تقریباً)
  صفر باشد گیر می‌کنند، هم در فلات‌ها و هم در کمینه‌ها؛ به مقدار شروع
  حساس‌اند.
\end{itemize}

\begin{quote}
\textbf{پیشرفته: چرا فلات‌ها روش‌های گرادیانی را به دام می‌اندازند.} اگر
\(\lambda\) خیلی بزرگ باشد، \(e^{-\lambda t} \approx 0\) در همه‌ی
\(t\)های مشاهده‌شده، پس هر دو مشتق جزئی بالا \(\approx 0\) هستند. روش
گرادیانی یک «کف صاف» می‌بیند و متوقف می‌شود و گاهی همگرایی موفق گزارش
می‌کند. \lr{Nelder\lr{--}Mead} مقادیر تابع را در نقاط دور مقایسه می‌کند و
می‌تواند از فلات بیرون بیاید. هر دو رفتار را در بخش \lr{D lab} می‌بینید.
\end{quote}

\textbf{خلاصه:}

{\def\LTcaptype{none} % do not increment counter
\begin{longtable}[]{@{}
  >{\raggedright\arraybackslash}p{(\linewidth - 6\tabcolsep) * \real{0.2500}}
  >{\raggedright\arraybackslash}p{(\linewidth - 6\tabcolsep) * \real{0.2500}}
  >{\raggedright\arraybackslash}p{(\linewidth - 6\tabcolsep) * \real{0.2500}}
  >{\raggedright\arraybackslash}p{(\linewidth - 6\tabcolsep) * \real{0.2500}}@{}}
\toprule\noalign{}
\begin{minipage}[b]{\linewidth}\raggedright
روش
\end{minipage} & \begin{minipage}[b]{\linewidth}\raggedright
نیاز به مشتق
\end{minipage} & \begin{minipage}[b]{\linewidth}\raggedright
کران
\end{minipage} & \begin{minipage}[b]{\linewidth}\raggedright
کاربرد معمول
\end{minipage} \\
\midrule\noalign{}
\endhead
\bottomrule\noalign{}
\endlastfoot
جستجوی شبکه‌ای & ندارد & ذاتی & رسم، مقدار شروع، \(k \le 2\) \\
\lr{Nelder\lr{--}Mead} & ندارد & دارد \lr{(scipy} \ensuremath{\geq} 1.7) & پیش‌فرض برای مدل‌های
شناختی کوچک؛ توابع هدف پرنویز \\
\lr{BFGS} & عددی یا تحلیلی & ندارد & توابع هدف هموار، پارامترهای بدون
قید \\
\lr{L-BFGS-B} & عددی یا تحلیلی & دارد & توابع هدف هموار با پارامترهای
کران‌دار \\
\end{longtable}
}

\paragraph{خواندن ریاضی}

گام گرادیانی
\(\boldsymbol\theta^{(m+1)} = \boldsymbol\theta^{(m)} - \eta\,\nabla \mathrm{SSE}(\boldsymbol\theta^{(m)})\)
این‌طور خوانده می‌شود: \emph{حدس جدید = حدس قبلی \ensuremath{-} اندازه‌ی گام × شیب}.
بالانویس \((m)\) شماره‌ی تکرار است، نه توان. \(\nabla \mathrm{SSE}\)
بردار مشتق‌های جزئی است (برای هر پارامتر یک شیب). علامت منفی گام را به
سمت سرازیری می‌برد.

فرمول مشتق
\(\partial\,\mathrm{SSE}/\partial\theta = -2\sum_i (y_i - \hat y_i)\,\partial \hat y_i/\partial\theta\)
از قاعده‌ی زنجیره‌ای می‌آید. می‌گوید: نقطه‌هایی که باقی‌مانده‌شان بزرگ است
\emph{و} پیش‌بینی در آن‌ها به \(\theta\) حساس است، \(\theta\) را قوی‌تر از
بقیه می‌کشند. در نقطه‌ی کمینه این کشش‌ها همدیگر را خنثی می‌کنند و جمع صفر
می‌شود.

\paragraph{برای یک مهندس}

\begin{itemize}
\tightlist
\item
  \textbf{کمترین مربعات} همان معیار تخمین کانال و برازش منحنی است؛ تفاوت
  اینجا این است که مدل‌ها نسبت به پارامترها \textbf{غیرخطی}‌اند
  (\(e^{-\lambda t}\))، پس جواب بسته‌ای از معادله‌های نرمال وجود ندارد و
  به جست‌وجوی تکراری نیاز داریم.
\item
  \textbf{کمترین مربعات وزنی} همان سفید کردن نویز \lr{(noise whitening)}
  است: تقسیم بر واریانس همان کاری است که وقتی نویز سفید نیست (نویز
  ناهمسان‌واریانس) پیش از \lr{LS} انجام می‌دهید.
\item
  \textbf{گرادیان نزولی با اندازه‌ی گام \(\eta\)} همان ایده‌ی \lr{steepest
  descent} در فیلتر تطبیقی \lr{LMS} است. \lr{BFGS} یک روش شبه‌نیوتنی است
  (تقریبی از \lr{Hessian} می‌سازد)، نزدیک به \lr{Gauss\lr{--}Newton} /
  \lr{Levenberg\lr{--}Marquardt} در \lr{LS} غیرخطی.
\item
  \textbf{یک دره‌ی دراز مورب} در سطح خطا یعنی مسئله \textbf{بدحالت
  \lr{(ill-conditioned)}} است: نسبت مقدار ویژه‌های \lr{Hessian} بزرگ است،
  مثل وقتی که رگرسورها به‌شدت همبسته‌اند. پارامترها در یک جهت خوب و در جهت
  دیگر بد تعیین می‌شوند.
\item
  \textbf{\lr{reparameterization}} (\(\lambda = e^{\eta}\)) یک تغییر
  متغیر است که مسئله‌ی مقید را به مسئله‌ی بدون قید تبدیل می‌کند و اغلب حالت
  \lr{(conditioning)} مسئله را هم بهتر می‌کند.
\item
  \textbf{تفاوت با کار مهندسی:} مجموعه داده‌ها کوچک‌اند (8 نقطه، هرکدام 50
  آیتم)، نویز زیاد است، و پارامترها باید از نظر روان‌شناختی \emph{تفسیر}
  شوند. پس کافی نیست منحنی پیش‌بینی‌شده خوب به نظر برسد؛ تک‌تک مقدارهای
  پارامتر باید قابل اعتماد باشند.
\end{itemize}

\subsubsection{4. مشکلات عملی و راه‌حل‌ها}

\textbf{1. مقادیر شروع.} از شروع‌های آگاهانه استفاده کنید: از روی نمودار
بخوانید، یک برازش خطی‌شده‌ی سریع (هفته‌ی 2) بگیرید، یا بهترین نقطه‌ی یک
شبکه‌ی درشت را بردارید.

\textbf{2. کمینه‌های محلی \ensuremath{\leftarrow} \lr{multi-start.}} بهینه‌ساز را از تعداد زیادی
نقطه‌ی شروع تصادفی اجرا کنید و بهترین جواب را نگه دارید. اگر شروع‌های
زیادی روی یک کمینه توافق کنند، می‌توانید نسبتاً مطمئن باشید که کمینه‌ی
سراسری است. \lr{multi-start} بیمه‌ی ارزانی است؛ برای مدل‌های بیش از دو
پارامتر همیشه از آن استفاده کنید.

\textbf{3. کران پارامترها.} بسیاری از پارامترها دامنه‌ی طبیعی دارند
(\(0 \le A \le 1\)، \(\lambda > 0\)). بدون قید، بهینه‌سازها با خیال راحت
\(A = 29.6\) برمی‌گردانند. سه رویکرد:

{\def\LTcaptype{none} % do not increment counter
\begin{longtable}[]{@{}
  >{\raggedright\arraybackslash}p{(\linewidth - 4\tabcolsep) * \real{0.3333}}
  >{\raggedright\arraybackslash}p{(\linewidth - 4\tabcolsep) * \real{0.3333}}
  >{\raggedright\arraybackslash}p{(\linewidth - 4\tabcolsep) * \real{0.3333}}@{}}
\toprule\noalign{}
\begin{minipage}[b]{\linewidth}\raggedright
رویکرد
\end{minipage} & \begin{minipage}[b]{\linewidth}\raggedright
چطور
\end{minipage} & \begin{minipage}[b]{\linewidth}\raggedright
توضیح
\end{minipage} \\
\midrule\noalign{}
\endhead
\bottomrule\noalign{}
\endlastfoot
بهینه‌ساز کران‌دار & \lr{\nolinkurl{method="L-BFGS-B"}} یا
\lr{\nolinkurl{"Nelder-Mead"}} با \lr{\nolinkurl{bounds=}} & ساده‌ترین؛
پیش‌فرض توصیه‌شده \\
\textbf{\lr{reparameterization}} & روی \(\eta\) بدون قید بهینه کنید و
تبدیل کنید: \(\lambda = e^{\eta}\) (مثبت)، \(A = 1/(1 + e^{-\alpha_A})\)
(بین 0 و 1) & با هر بهینه‌سازی کار می‌کند؛ در مدل‌سازی بیزی هم استفاده
می‌شود (هفته‌ی 13) \\
جریمه & وقتی پارامتر خارج از دامنه است یک \lr{SSE} خیلی بزرگ برگردانید &
ناپیوستگی می‌سازد؛ از آن پرهیز کنید \\
\end{longtable}
}

\textbf{4. \lr{non-identifiability.}} گاهی \emph{چند} ترکیب پارامتر به
یک اندازه خوب برازش می‌شوند. مثالش مدلی دوجزئی است که برچسب دو جزئش قابل
جابه‌جایی است. این ویژگی خود مدل است، نه شکست بهینه‌ساز، و
\lr{multi-start} آن را آشکار می‌کند (بخش \lr{D lab).}

\textbf{5. اعتماد کن، ولی بررسی کن.}
\lr{\nolinkurl{result.success == True}} فقط یعنی قاعده‌ی توقف الگوریتم
برقرار شده است. همیشه (1) منحنی برازش‌شده را روی داده رسم کنید، (2) به
باقی‌مانده‌ها نگاه کنید، و (3) بررسی کنید برآوردها در دامنه‌ی معقول باشند.

\textbf{چک‌لیست برازش (از این به بعد در هر \lr{lab} از آن استفاده کنید):}

\begin{enumerate}
\def\labelenumi{\arabic{enumi}.}
\tightlist
\item
  داده‌ها را رسم کنید.
\item
  مدل و تابع ناهمخوانی را بنویسید؛ آن‌ها را روی چند مقدار پارامتر دلخواه
  آزمون کنید.
\item
  کران‌ها و مقادیر شروع معقول انتخاب کنید.
\item
  بهینه کنید، ترجیحاً با \lr{multi-start.}
\item
  پیام همگرایی، نمودار، باقی‌مانده‌ها و دامنه‌ی پارامترها را بررسی کنید.
\item
  یک بررسی بازیابی انجام دهید: از پارامترهای برازش‌شده داده شبیه‌سازی
  کنید، دوباره برازش کنید و مقایسه کنید (هفته‌ی 7 این را نظام‌مند می‌کند).
\end{enumerate}

\begin{quote}
\textbf{پیشرفته: پارامترهای شرطاً خطی.} در \(f(t) = A\,e^{-\lambda t}\)،
برای هر \(\lambda\) ثابت بهترین \(A\) فرم بسته دارد، چون مدل نسبت به
\(A\) خطی است: \(\hat A(\lambda) = \sum_i y_i g_i / \sum_i g_i^2\) با
\(g_i = e^{-\lambda t_i}\). جای‌گذاری آن مسئله را به یک مسئله‌ی یک‌بعدی فقط
بر حسب \(\lambda\) تبدیل می‌کند. این ترفند (حذف پارامترهای خطی یا
\emph{\lr{variable projection}}) برازش را سریع‌تر و پایدارتر می‌کند، و
تابع یک‌بعدی حاصل با \emph{\lr{profile likelihood}} هفته‌ی 5 مرتبط است.
\end{quote}

\paragraph{برای مدرس}

\begin{itemize}
\tightlist
\item
  \textbf{اول منظره را نشان دهید.} نمودار تراز \lr{SSE} روی
  \((A, \lambda)\) (بخش \lr{B lab)} همه‌ی ایده‌های بعدی را ملموس می‌کند:
  کمینه، دره‌ی مورب، فلات. آن را زود روی پروژکتور بیاورید و مدام به آن
  برگردید.
\item
  \textbf{\lr{Nelder\lr{--}Mead} را اجرا کنید.} یک مثلث روی نمودار تراز بکشید
  و یک بازتاب \lr{(reflection)} را با دست انجام دهید (قدم 2 مثال حل‌شده).
  دانشجوها «مثلثی که قل می‌خورد» را خیلی بهتر از جدول حرکت‌ها به خاطر
  می‌سپارند.
\item
  \textbf{نقطه‌ی قوت شما:} شما \lr{gradient descent} و کمترین مربعات را
  می‌شناسید. وقت کمتری روی الگوریتم‌ها و وقت بیشتری روی \emph{حالت‌های
  شکست} بگذارید (بخش \lr{D lab):} بهینه‌سازی که در \(A = 29.6\) اعلام
  موفقیت می‌کند، به‌یادماندنی‌ترین درس این هفته است.
\item
  \textbf{سؤال‌های رایج.}

  \begin{itemize}
  \tightlist
  \item
    \emph{«از کدام بهینه‌ساز استفاده کنم؟»} برای مدل‌های شناختی کوچک:
    \lr{Nelder\lr{--}Mead} یا \lr{L-BFGS-B} با کران، همیشه همراه با
    \lr{multi-start.}
  \item
    \emph{«چرا همیشه از \lr{grid} استفاده نکنیم؟»} هزینه با تعداد
    پارامترها به‌طور نمایی زیاد می‌شود.
  \item
    \emph{«آیا \lr{SSE} معیار درستی است؟»} فقط اگر نویز گاوسی و با
    واریانس برابر باشد؛ هفته‌ی بعد درست‌نمایی جای آن را می‌گیرد.
  \end{itemize}
\end{itemize}

\begin{center}\rule{0.5\linewidth}{0.5pt}\end{center}

\subsection{مثال حل‌شده}

\textbf{داده.} یک شرکت‌کننده نسبت‌های \(y = (0.62, 0.41, 0.19)\) را در
\(t = (1, 3, 7)\) روز به یاد می‌آورد. مدل: \(f(t) = A\,e^{-\lambda t}\).

\subsubsection{قدم 1: یک جستجوی شبکه‌ای کوچک با دست}

برای \((A, \lambda) = (0.8, 0.2)\): پیش‌بینی‌ها \(0.8e^{-0.2} = 0.655\)،
\(0.8e^{-0.6} = 0.439\)، \(0.8e^{-1.4} = 0.197\)، پس

\[
\mathrm{SSE} = (0.62 - 0.655)^2 + (0.41 - 0.439)^2 + (0.19 - 0.197)^2 = 0.00123 + 0.00084 + 0.00005 \approx 0.00212.
\]

با تکرار این کار برای یک شبکه‌ی 2 × 3:

{\def\LTcaptype{none} % do not increment counter
\begin{longtable}[]{@{}llll@{}}
\toprule\noalign{}
& \(\lambda = 0.15\) & \(\lambda = 0.20\) & \(\lambda = 0.25\) \\
\midrule\noalign{}
\endhead
\bottomrule\noalign{}
\endlastfoot
\(A = 0.8\) & 0.02281 & \textbf{0.00212} & 0.00364 \\
\(A = 0.9\) & 0.06638 & 0.02172 & 0.00791 \\
\end{longtable}
}

بهترین نقطه‌ی شبکه \((0.8, 0.20)\) است. به کجی دقت کنید: \((0.9, 0.25)\)
خیلی بهتر از \((0.9, 0.15)\) است. \(A\) بزرگ‌تر یک \(\lambda\) بزرگ‌تر
«می‌خواهد»؛ این همان دره‌ی \lr{trade-off} است.

\subsubsection{\texorpdfstring{قدم 2: یک تکرار \lr{Nelder\lr{--}Mead} با
دست}{قدم 2: یک تکرار Nelder\lr{--}Mead با دست}}

با سیمپلکس (مثلث) سه نقطه‌ی شبکه شروع کنید:

{\def\LTcaptype{none} % do not increment counter
\begin{longtable}[]{@{}llll@{}}
\toprule\noalign{}
رأس & \((A, \lambda)\) & \lr{SSE} & رتبه \\
\midrule\noalign{}
\endhead
\bottomrule\noalign{}
\endlastfoot
\(\boldsymbol\theta_b\) & (0.8, 0.25) & 0.00364 & بهترین \\
\(\boldsymbol\theta_g\) & (0.8, 0.15) & 0.02281 & دومین بدترین \\
\(\boldsymbol\theta_w\) & (0.9, 0.15) & 0.06638 & بدترین \\
\end{longtable}
}

\begin{enumerate}
\def\labelenumi{\arabic{enumi}.}
\tightlist
\item
  \textbf{مرکز ثقل} همه‌ی رأس‌ها جز بدترین:
  \(\bar{\boldsymbol\theta} = \big((0.8 + 0.8)/2,\ (0.25 + 0.15)/2\big) = (0.8,\allowbreak 0.20)\).
\item
  \textbf{بازتاب} بدترین رأس از میان مرکز ثقل:
  \(\boldsymbol\theta_r = (0.8,\allowbreak 0.20) + \big((0.8,\allowbreak 0.20) - (0.9,\allowbreak 0.15)\big) = (0.7,\allowbreak 0.25)\).
\item
  \textbf{ارزیابی:} پیش‌بینی‌ها \((0.545, 0.331, 0.122)\)، پس
  \(\mathrm{SSE}(\boldsymbol\theta_r) = 0.01657\).
\item
  \textbf{تصمیم:} \(0.00364 \le 0.01657 < 0.02281\). نقطه‌ی بازتاب‌شده از
  دومین بدترین بهتر است ولی از بهترین بهتر نیست، پس \textbf{بازتاب را
  می‌پذیریم} (بدون گسترش).
\item
  \textbf{سیمپلکس جدید:}
  \(\{(0.8,\allowbreak 0.25),\ (0.7,\allowbreak 0.25),\ (0.8,\allowbreak 0.15)\}\)، و بدترین \lr{SSE} از
  0.066 به 0.023 کم شده است.
\end{enumerate}

\subsubsection{قدم 3: بگذارید کامپیوتر تمامش کند}

اجرای
\lr{\nolinkurl{scipy.optimize.minimize(..., method="Nelder-Mead")}} از
\((0.8, 0.2)\) به \(\hat A = 0.754\)، \(\hat\lambda = 0.200\) با
\(\mathrm{SSE} = 3.7 \times 10^{-5}\) \lr{(RMSD} \(= 0.0035\)) همگرا
می‌شود.

\textbf{تفسیر.} مدل برازش‌شده می‌گوید این شرکت‌کننده بلافاصله بعد از
یادگیری حدود 75\% را به یاد می‌آورد و هر روز حدود 18\% از باقی‌مانده را از
دست می‌دهد (\(1 - e^{-0.2} = 0.181\))، با نیمه‌عمر
\(\ln 2 / 0.2 \approx 3.5\) روز. با فقط سه داده و دو پارامتر، برازش
تقریباً کامل چیز زیادی درباره‌ی \emph{درست بودن} مدل نمایی نمی‌گوید. هر
منحنی نزولی دوپارامتری معقولی سه نقطه را خوب برازش می‌کند (هفته‌ی 6).

\begin{center}\rule{0.5\linewidth}{0.5pt}\end{center}

\subsection{\texorpdfstring{\lr{lab} (2 ساعت)}{lab (2 ساعت)}}

\subsubsection{اهداف}

\begin{enumerate}
\def\labelenumi{\arabic{enumi}.}
\tightlist
\item
  نوشتن یک تابع \lr{SSE} و بررسی سطح خطای دو مدل فراموشی.
\item
  برازش با جستجوی شبکه‌ای، سپس با \lr{Nelder\lr{--}Mead} و \lr{L-BFGS-B}؛
  مقایسه با برازش‌های خطی‌شده‌ی هفته‌ی قبل.
\item
  ایجاد و رفع شکست‌های برازش: فلات‌ها، کمینه‌های محلی و برآوردهای خارج از
  دامنه.
\item
  \textbf{ارزیابی} برازش‌ها \lr{(RMSD}، باقی‌مانده‌ها، مقایسه با مقادیر
  واقعی) و \textbf{تفسیر} پارامترها.
\end{enumerate}

\textbf{زمان‌بندی:} بخش \lr{A} 15 دقیقه · بخش \lr{B} 25 دقیقه · بخش
\lr{C} 25 دقیقه · بخش \lr{D} 35 دقیقه · بخش \lr{E} 20 دقیقه.

\subsubsection{آماده‌سازی و داده}

\begin{quote}
\textbf{\lr{Code} · کد:} \lr{\nolinkurl{code/week03.py}}, \lr{cell} 1 ·
سلول 1
\end{quote}

\subsubsection{\texorpdfstring{بخش \lr{A:} آشنایی با تابع ناهمخوانی (15
دقیقه)}{بخش A: آشنایی با تابع ناهمخوانی (15 دقیقه)}}

\begin{quote}
\textbf{\lr{Code} · کد:} \lr{\nolinkurl{code/week03.py}}, \lr{cell} 2 ·
سلول 2
\end{quote}

\textbf{خروجی مورد انتظار:}

\Needspace{9\baselineskip}
\begin{latin}
\begin{Verbatim}
theta = (0.9, 0.5):  power SSE = 0.0349   exponential SSE = 0.4511
theta = (0.8, 0.4):  power SSE = 0.0328   exponential SSE = 0.4145
theta = (0.9, 0.1):  power SSE = 0.9047   exponential SSE = 0.2192
theta = (0.6, 0.05):  power SSE = 0.3303   exponential SSE = 0.1145
\end{Verbatim}
\end{latin}

پیش از اجرای بخش‌های بعد \textbf{پیش‌بینی کنید}: کدام مقادیر پارامتر
کمترین \lr{SSE} را برای هر مدل می‌دهند؟ حدستان را یادداشت کنید.

\subsubsection{\texorpdfstring{بخش \lr{B:} جستجوی شبکه‌ای و سطح خطا (25
دقیقه)}{بخش B: جستجوی شبکه‌ای و سطح خطا (25 دقیقه)}}

\begin{quote}
\textbf{\lr{Code} · کد:} \lr{\nolinkurl{code/week03.py}}, \lr{cell} 3 ·
سلول 3
\end{quote}

\textbf{خروجی مورد انتظار:}

\Needspace{7\baselineskip}
\begin{latin}
\begin{Verbatim}
exp: grid minimum A = 0.655, rate = 0.0560, SSE = 0.10607  (42300 evaluations)
pow: grid minimum A = 0.870, rate = 0.4239, SSE = 0.02592  (42300 evaluations)
\end{Verbatim}
\end{latin}

\textbf{به این دقت کنید:} هر دو سطح یک حوضه دارند، ولی یک \textbf{دره‌ی
کشیده و مورب}: \(A\) بالاتر با زوال بیشتر همراه است. این همان
\lr{trade-off} پارامترهاست. برای مدل توانی، مقدار واقعی (× سفید) داخل
دره است ولی در ته آن نیست: با داده‌ی محدود، بهترین برازش هیچ‌وقت دقیقاً
همان مقدار واقعی نیست.

\subsubsection{\texorpdfstring{بخش \lr{C:} برازش با یک بهینه‌ساز (25
دقیقه)}{بخش C: برازش با یک بهینه‌ساز (25 دقیقه)}}

\begin{quote}
\textbf{\lr{Code} · کد:} \lr{\nolinkurl{code/week03.py}}, \lr{cell} 4 ·
سلول 4
\end{quote}

\textbf{خروجی مورد انتظار:}

\Needspace{14\baselineskip}
\begin{latin}
\begin{Verbatim}
model            method      A   rate    SSE   RMSD  f_evals success
  exp          L-BFGS-B 0.6568 0.0566 0.1061 0.1151     39.0    True
  exp       Nelder-Mead 0.6568 0.0566 0.1061 0.1151     69.0    True
  exp              grid 0.6550 0.0560 0.1061 0.1151  42300.0    True
  exp linearized (Wk 2) 0.5500 0.0320 0.1373 0.1310      NaN     NaN
  pow          L-BFGS-B 0.8712 0.4251 0.0259 0.0569     36.0    True
  pow       Nelder-Mead 0.8711 0.4251 0.0259 0.0569     71.0    True
  pow              grid 0.8700 0.4239 0.0259 0.0569  42300.0    True
  pow linearized (Wk 2) 0.7970 0.3700 0.0315 0.0628      NaN     NaN
\end{Verbatim}
\end{latin}

\begin{quote}
\textbf{\lr{Code} · کد:} \lr{\nolinkurl{code/week03.py}}, \lr{cell} 5 ·
سلول 5
\end{quote}

\textbf{به این‌ها دقت کنید:} - \lr{Nelder\lr{--}Mead} و \lr{L-BFGS-B} تا سه یا
چهار رقم اعشار با هم توافق دارند؛ شبکه در حد فاصله‌ی نقاطش (0.005 برای
\(A\)) با آن‌ها فرق دارد. بهینه‌سازها فقط 36 تا 71 ارزیابی تابع لازم
داشتند، به‌جای 42 300. - کمترین مربعات \textbf{\lr{SSE} کمتری} از
برازش‌های خطی‌شده‌ی هفته‌ی 2 می‌دهد (دقیقاً همین کمیت را کمینه می‌کند)، و
برآوردهای مدل توانی خیلی به مقدار واقعی نزدیک‌ترند. - مسیر اول سریع وارد
دره می‌شود و بعد آرام در امتدادش می‌خزد؛ یک الگوی معمول.

\subsubsection{\texorpdfstring{بخش \lr{D:} وقتی برازش خراب می‌شود (35
دقیقه)}{بخش D: وقتی برازش خراب می‌شود (35 دقیقه)}}

\textbf{\lr{D1.} یک فلات و یک «موفقیت» کاذب.} برازش نمایی را با یک نرخ
زوال بی‌معنی، \(\lambda = 20\) در روز، شروع کنید:

\begin{quote}
\textbf{\lr{Code} · کد:} \lr{\nolinkurl{code/week03.py}}, \lr{cell} 6 ·
سلول 6
\end{quote}

\textbf{خروجی مورد انتظار:}

\Needspace{8\baselineskip}
\begin{latin}
\begin{Verbatim}
BFGS         A =   29.571  lambda =  14.244  SSE = 1.1616  success = True  (Optimization terminated successfully.)
L-BFGS-B     A =   29.905  lambda =  14.289  SSE = 1.1616  success = True  (CONVERGENCE: NORM_OF_PROJECTED_GRADIENT_<=_PGTOL)
Nelder-Mead  A =    0.657  lambda =   0.057  SSE = 0.1061  success = True  (Optimization terminated successfully.)
\end{Verbatim}
\end{latin}

(متن پیام \lr{L-BFGS-B} بین نسخه‌های \lr{SciPy} کمی فرق می‌کند.) دو روش
گرادیانی \textbf{موفقیت} گزارش می‌کنند ولی برازشی بی‌معنی برمی‌گردانند:
\(A\) حدود 30 است (احتمال یادآوری 3000\%!) و منحنی پیش‌بینی‌شده عملاً در
همه‌ی فاصله‌ها صفر است. آن‌ها روی فلاتی شروع کردند که گرادیان در آن تقریباً
صفر است (کادر «پیشرفته» بخش 3.3 جزوه). \lr{Nelder\lr{--}Mead} فرار کرد.
درس‌ها: (1) همیشه برآوردها و نمودارها را بررسی کنید، نه فقط
\lr{\nolinkurl{success}} را؛ (2) از کران استفاده کنید؛ (3) از مقادیر
شروع معقول یا متعدد استفاده کنید.

\textbf{\lr{D2. reparameterization.}} پارامترها را بدون قید کنید تا
\emph{هر} بهینه‌سازی دامنه‌های \(0 < A < 1\) و \(\lambda > 0\) را رعایت
کند:

\begin{quote}
\textbf{\lr{Code} · کد:} \lr{\nolinkurl{code/week03.py}}, \lr{cell} 7 ·
سلول 7
\end{quote}

\textbf{خروجی مورد انتظار:}

\Needspace{6\baselineskip}
\begin{latin}
\begin{Verbatim}
BFGS on reparameterized power model: A = 0.8712, d = 0.4251, SSE = 0.02592
\end{Verbatim}
\end{latin}

همان جواب بخش \lr{C}، و تضمین می‌شود که برآوردها در دامنه باشند.

\textbf{\lr{D3.} کمینه‌های محلی و \lr{multi-start.}} یک ایده‌ی رایج در
پژوهش‌های حافظه این است که فراموشی بازتاب دو رد حافظه \lr{(memory trace)}
است: یک جزء \textbf{سریع} که زود زوال پیدا می‌کند و یک جزء \textbf{کند}
که می‌ماند:

\[
f(t) = A_1 e^{-\lambda_1 t} + A_2 e^{-\lambda_2 t}.
\]

این مدل \(k = 4\) پارامتر دارد و سطح خطایش چند حوضه دارد. آن را از 50
نقطه‌ی شروع تصادفی برازش کنید:

\begin{quote}
\textbf{\lr{Code} · کد:} \lr{\nolinkurl{code/week03.py}}, \lr{cell} 8 ·
سلول 8
\end{quote}

\textbf{خروجی مورد انتظار:}

\Needspace{19\baselineskip}
\begin{latin}
\begin{Verbatim}
Distinct final SSE values (rounded) and how often each was reached:
  SSE = 0.0201: 42 starts
  SSE = 0.0203: 1 starts
  SSE = 0.0337: 1 starts
  SSE = 0.0340: 1 starts
  SSE = 0.0359: 1 starts
  SSE = 0.1064: 1 starts
  SSE = 0.2010: 1 starts
  SSE = 0.2456: 1 starts
  SSE = 0.9092: 1 starts
Best 3 solutions (A1, lam1, A2, lam2):
   [0.58  0.36  0.264 0.   ] SSE = 0.02013
   [0.264 0.    0.58  0.36 ] SSE = 0.02013
   [0.264 0.    0.58  0.36 ] SSE = 0.02013
\end{Verbatim}
\end{latin}

\textbf{به این‌ها دقت کنید:} - حدود 40 شروع از 50 در همان بهترین حوضه
تمام می‌شوند، ولی بقیه (حدود 1 از 6) در \textbf{کمینه‌های محلی بدتر} گیر
می‌کنند؛ بعضی با \lr{SSE} چند برابر بزرگ‌تر. (تعداد دقیق بین نسخه‌های
\lr{SciPy} کمی فرق دارد؛ با \lr{SciPy} 1.18 برابر 41 از 50 است.) یک
اجرای تنها از یک شروع بدشانس جواب غلطی می‌داد، بدون هیچ هشداری. - در میان
بهترین جواب‌ها، اجزای سریع و کند ممکن است \textbf{با برچسب‌های جابه‌جا}
ظاهر شوند (\(A_1,\allowbreak \lambda_1 \leftrightarrow A_2,\allowbreak \lambda_2\)). هر دو
همان منحنی را توصیف می‌کنند. این یک \lr{non-identifiability} ذاتی مدل
است، و راه‌حل رایج اعمال قید ترتیبی مثل \(\lambda_1 > \lambda_2\) است. -
در بهترین جواب، یکی از اجزا نرخ دقیقاً 0 دارد (یک جزء ثابت و بدون زوال به
اندازه‌ی 0.264). پس مدل دوردی برازش‌شده به مدل نمایی با مجانب تمرین 4
تبدیل شده، با همان \lr{SSE} (0.0201). آیا این یک «حافظه‌ی دائمی» معنادار
است یا فقط پارامترهای اضافه نویز را جذب کرده‌اند؟ یادتان باشد داده‌ها از
یک مدل \emph{توانی} آمده‌اند. برازش به‌تنهایی نمی‌تواند جواب دهد (هفته‌های
6\lr{--}7).

\subsubsection{\texorpdfstring{بخش \lr{E:} ارزیابی و تفسیر (20
دقیقه)}{بخش E: ارزیابی و تفسیر (20 دقیقه)}}

\begin{quote}
\textbf{\lr{Code} · کد:} \lr{\nolinkurl{code/week03.py}}, \lr{cell} 9 ·
سلول 9
\end{quote}

\textbf{خروجی مورد انتظار:}

\Needspace{8\baselineskip}
\begin{latin}
\begin{Verbatim}
Power model: relative loss rate at t =  1 days: 0.213 per day
Power model: relative loss rate at t = 20 days: 0.020 per day
Exponential model: constant relative loss rate 0.057 per day
\end{Verbatim}
\end{latin}

\textbf{تفسیر (برای هرکدام 2 تا 3 جمله در \lr{notebook} بنویسید):}

\begin{enumerate}
\def\labelenumi{\arabic{enumi}.}
\tightlist
\item
  الگوهای باقی‌مانده‌ها را مقایسه کنید. کدام مدل الگوی نظام‌مند نشان می‌دهد،
  و آن الگو درباره‌ی شکل فراموشی در این داده‌ها چه می‌گوید؟
\item
  نرخ نسبی فراموشی مدل توانی از حدود 0.21 در روز بعد از 1 روز به حدود
  0.02 در روز بعد از 20 روز می‌رسد. به زبان ساده توضیح دهید این چه ادعایی
  درباره‌ی حافظه دارد.
\item
  برآوردهای کمترین مربعات (\(\hat A \approx 0.87\)،
  \(\hat d \approx 0.43\)) از برآوردهای خطی‌شده‌ی هفته‌ی قبل به مقدار واقعی
  نزدیک‌ترند ولی هنوز با آن برابر نیستند. آیا این شکست روش است؟
  (راهنمایی: اگر آزمایش را با یک نمونه‌ی تصادفی جدید تکرار می‌کردید چه
  می‌شد؟)
\end{enumerate}

\subsubsection{تمرین‌ها (به ترتیب سختی)}

\textbf{تمرین 1 \ensuremath{\star} (کمینه‌کننده‌ی یکسان).} مدل توانی را یک بار با کمینه
کردن \lr{SSE}، یک بار \lr{MSE} و یک بار \lr{RMSD} برازش کنید. نشان دهید
برآوردها (در حد تلرانس بهینه‌ساز) یکسان‌اند. چرا باز هم ترجیح می‌دهیم
\lr{RMSD} را \emph{گزارش} کنیم؟

\textbf{تمرین 2 \ensuremath{\star} (ابینگهاوس با کمترین مربعات).} داده‌های \lr{savings}
ابینگهاوس از تمرین 3 هفته‌ی 2 (مقادیر با علامت {[}بررسی شود{]}) را برای
هر دو مدل با کمترین مربعات دوباره برازش کنید. پارامترها و \lr{RMSD} را
با برازش‌های خطی‌شده‌ی هفته‌ی 2 مقایسه کنید.

\textbf{تمرین 3 \ensuremath{\star}\ensuremath{\star} (کمترین مربعات وزنی).} نسبت‌های دوجمله‌ای واریانس
\(f(1-f)/n\) دارند. \lr{WSSE} (بخش 1 جزوه) را با استفاده از \(f\)
\textbf{پیش‌بینی‌شده} در واریانس پیاده‌سازی کنید، مدل توانی را دوباره برازش
کنید و برآوردها را با کمترین مربعات ساده مقایسه کنید. کدام نقاط داده
تأثیر بیشتری پیدا می‌کنند و کدام کمتر؟ (هفته‌ی بعد می‌بینید که بیشینه
درست‌نمایی این وزن‌دهی را خودکار انجام می‌دهد.)

\textbf{تمرین 4 \ensuremath{\star}\ensuremath{\star} (پارامتر بیشتر، برازش بهتر؟).} مدل سه‌پارامتری
\textbf{نمایی با مجانب} \(f(t) = B + (A - B)e^{-\lambda t}\) را که در آن
\(B\) مجانب بلندمدت است، با کران‌های \(0 \le B \le A \le 1\) برازش کنید
(راهنمایی: \lr{reparameterize} کنید، یا قید را بعد از برازش بررسی کنید).
\emph{انتظار:} \lr{SSE} \ensuremath{\approx} 0.020 که \textbf{کمتر} از \lr{SSE} \ensuremath{\approx} 0.026 مدل
توانی است، با اینکه داده‌ها از مدل توانی تولید شده‌اند! این درباره‌ی
استفاده از \lr{SSE} برای انتخاب بین مدل‌هایی با تعداد پارامتر متفاوت چه
می‌گوید؟ (در هفته‌های 6\lr{--}7 جواب درست این را می‌دهیم.)

\textbf{تمرین 5 \ensuremath{\star}\ensuremath{\star}\ensuremath{\star} (پیشرفته: \lr{Nelder\lr{--}Mead} از صفر).}
\lr{Nelder\lr{--}Mead} را برای \(k = 2\) با حرکت‌های بخش 3.2 جزوه (بازتاب،
گسترش، انقباض داخلی/خارجی، جمع‌شدن) و یک قاعده‌ی توقف بر اساس پراکندگی
مقادیر \lr{SSE} در سیمپلکس پیاده‌سازی کنید. آن را روی مدل توانی از
\((0.8, 0.2)\) اجرا کنید، سیمپلکس را هر چند تکرار یک بار روی سطح خطا رسم
کنید، و نتیجه و تعداد ارزیابی‌های تابع را با \lr{\nolinkurl{scipy}}
مقایسه کنید.

\subsubsection{\texorpdfstring{\lr{lab check-off}}{lab check-off}}

به مدرس یا دستیار آموزشی نشان دهید: (1) نمودارهای سطح خطا با کمینه‌های
شبکه (بخش \lr{B)}، (2) جدول مقایسه از بخش \lr{C}، (3) هیستوگرام
\lr{multi-start} (بخش \lr{D3)} با یک جمله درباره‌ی آنچه نشان می‌دهد، و (4)
جوابتان به سؤال تفسیری 1 (بخش \lr{E).}

\begin{center}\rule{0.5\linewidth}{0.5pt}\end{center}

\subsection{\texorpdfstring{ارتباط با هوش مصنوعی و
\lr{LLM}ها}{ارتباط با هوش مصنوعی و LLMها}}

\textbf{قانون‌های مقیاس \lr{(scaling laws)} همان برازش منحنی‌اند.} یکی از
تأثیرگذارترین نتایج هوش مصنوعی امروز این است که خطای آزمون \lr{(test
loss)} مدل‌های زبانی با اندازه‌ی مدل \(N\)، اندازه‌ی داده و محاسبات به‌طور
هموار کم می‌شود و از یک \textbf{قانون توانی} با مجانب پیروی می‌کند:
\(L(N) = E + B\,N^{-\alpha}\) \lr{(Kaplan et al.,} 2020; \lr{Hoffmann et
al.,} 2022). شرکت‌ها از چنین برازش‌هایی استفاده می‌کنند تا تصمیم بگیرند مدل
بعدی چقدر بزرگ باشد. از نظر ریاضی، این همان مسئله‌ی برازش منحنی فراموشی
است: یک مدل غیرخطی با سه پارامتر که با کمترین مربعات، یک بهینه‌ساز عددی و
چند نقطه‌ی شروع برازش می‌شود. دام‌هایش هم همان‌اند: برون‌یابی
\lr{(extrapolation)} خیلی دورتر از داده، و \lr{trade-off} پارامترها
(اینجا \(E\) و \(\alpha\)).

\begin{quote}
\textbf{\lr{Code} · کد:} \lr{\nolinkurl{code/week03.py}}, \lr{cell} 10 ·
سلول 10
\end{quote}

\textbf{خروجی مورد انتظار:}

\Needspace{7\baselineskip}
\begin{latin}
\begin{Verbatim}
E = 1.70, B = 3.22, alpha = 0.358, SSE = 0.00424
Extrapolated loss at N = 1e11: 1.754
\end{Verbatim}
\end{latin}

\textbf{امتحانش کنید (بدون نیاز به \lr{API).}} - \ensuremath{\star} یک قانون توانی خالص
بدون مجانب (\(E = 0\)) برازش کنید. \lr{SSE} و برون‌یابی تا
\(N = 10^{11}\) را مقایسه کنید. کدام نتیجه‌گیری درباره‌ی «بزرگ‌تر بهتر است»
عوض می‌شود؟ - \ensuremath{\star}\ensuremath{\star} فقط چهار مدل کوچک‌تر را برازش کنید و تا چهار مدل بزرگ‌تر
برون‌یابی کنید. پیش‌بینی چقدر خطا دارد؟ با \lr{seed}های مختلف تکرار کنید
تا ببینید برون‌یابی چقدر پایدار است. - \ensuremath{\star}\ensuremath{\star} سطح \lr{SSE} را روی \(E\) و
\(\alpha\) رسم کنید (با \(B\) در بهترین مقدارش برای هر جفت). آیا مثل
\lr{lab} هفته‌ی 3 یک \lr{ridge} وجود دارد؟

\emph{مطالعه‌ی اختیاری:} \lr{Kaplan, J., et al.~(2020). Scaling laws for
neural language models.} \emph{\lr{arXiv:2001.08361}}. \lr{Hoffmann, J.,
et al.~(2022). Training compute-optimal large language models.}
\emph{\lr{arXiv:2203.15556}}.

\begin{center}\rule{0.5\linewidth}{0.5pt}\end{center}

\subsection{اشتباهات رایج و برداشت‌های نادرست}

\begin{enumerate}
\def\labelenumi{\arabic{enumi}.}
\tightlist
\item
  \textbf{«\lr{\nolinkurl{success=True}} یعنی برازش خوب است.»} فقط یعنی
  یک قاعده‌ی توقف برقرار شده است (بخش \lr{D1).} همیشه برآوردها، نمودارها
  و باقی‌مانده‌ها را وارسی کنید.
\item
  \textbf{یک شروع کافی است.} برای هر مدلی فراتر از مدل‌های ساده‌ی
  دوپارامتری \lr{multi-start} اجرا کنید. هزینه‌اش کم است؛ خطر گزارش یک
  کمینه‌ی محلی کم نیست.
\item
  \textbf{نادیده گرفتن دامنه‌ی پارامترها.} بهینه‌سازهای بدون کران می‌توانند
  مقادیر ناممکن برگردانند (نرخ منفی، احتمال بیشتر از 1). از کران یا
  \lr{reparameterization} استفاده کنید.
\item
  \textbf{مقایسه‌ی \lr{SSE} بین مدل‌هایی با تعداد پارامتر متفاوت.} مدل
  منعطف‌تر تقریباً همیشه در \lr{SSE} برنده است (تمرین 4). این سؤالی
  درباره‌ی \emph{ارزیابی} مدل است، نه برآورد (هفته‌های 6\lr{--}7).
\item
  \textbf{مقایسه‌ی \lr{SSE} بین مجموعه‌داده‌هایی با اندازه‌ی متفاوت.}
  \lr{SSE} با \(N\) بزرگ می‌شود؛ برای مقایسه از \lr{RMSD} استفاده کنید.
\item
  \textbf{کمینه‌کننده را همان مقدار واقعی دانستن.}
  \(\hat{\boldsymbol\theta}\) برآوردی با خطای نمونه‌گیری است. نمونه‌ی دیگر
  برآورد دیگری می‌دهد (هفته‌ی 5 این را کمّی می‌کند).
\item
  \textbf{فراموش کردن فرض کمترین مربعات.} \lr{SSE} ساده نویز گاوسی و
  برابر برای همه‌ی مشاهدات فرض می‌کند. نسبت‌ها، شمارش‌ها و زمان‌های واکنش
  معمولاً این فرض را نقض می‌کنند.
\item
  \textbf{استفاده از جستجوی شبکه‌ای برای تعداد پارامتر زیاد.} هزینه منفجر
  می‌شود (\(100^k\) ارزیابی). شبکه را فقط برای رسم و مقادیر شروع به کار
  ببرید.
\end{enumerate}

\begin{center}\rule{0.5\linewidth}{0.5pt}\end{center}

\subsection{خودآزمایی}

\textbf{1.} چرا \lr{SSE}، \lr{MSE} و \lr{RMSD} کمینه‌کننده‌ی یکسانی دارند؟
کدام را گزارش می‌کنید و چرا؟ \textbf{پاسخ} \lr{MSE} = \lr{SSE/N} و
\lr{RMSD} = \lr{\ensuremath{\surd}MSE} تبدیل‌های اکیداً صعودی \lr{SSE} هستند \lr{(N} ثابت
است)، پس همان پارامترها هر سه را کمینه می‌کنند. \lr{RMSD} را گزارش کنید،
چون بر حسب واحد داده است و بین مجموعه‌داده‌هایی با اندازه‌ی متفاوت قابل
مقایسه است.

\textbf{2.} \lr{trade-off} پارامترها چیست و روی سطح خطا چگونه دیده
می‌شود؟ \textbf{پاسخ} تغییر یک پارامتر را می‌شود با تغییر پارامتر دیگری
جبران کرد، طوری که پیش‌بینی‌ها (و \lr{SSE)} تقریباً عوض نشوند. روی سطح به
شکل یک دره‌ی کشیده و مورب با \lr{SSE} پایین دیده می‌شود. در این حالت
برآوردها همبسته‌اند و در جهت دره ضعیف تعیین می‌شوند.

\textbf{3.} با 50 مقدار شبکه برای هر پارامتر، جستجوی شبکه‌ای برای یک مدل
6 پارامتری به چند ارزیابی \lr{SSE} نیاز دارد؟ این یعنی چه؟ \textbf{پاسخ}
\(50^6 \approx 1.56 \times 10^{10}\) ارزیابی. جستجوی شبکه‌ای فراتر از 2
تا 3 پارامتر عملی نیست (نفرین ابعاد)؛ به بهینه‌سازهای تکراری نیاز داریم.

\textbf{4.} در یک تکرار \lr{Nelder\lr{--}Mead}، نقطه‌ی بازتاب‌شده از بهترین رأس
فعلی بهتر است. الگوریتم بعد چه می‌کند؟ \textbf{پاسخ} یک گسترش امتحان
می‌کند: نقطه‌ای در همان جهت با فاصله‌ی دو برابر از مرکز ثقل. اگر نقطه‌ی
گسترش‌یافته از نقطه‌ی بازتاب‌شده بهتر باشد آن را نگه می‌دارد، وگرنه نقطه‌ی
بازتاب‌شده را.

\textbf{5.} چرا یک بهینه‌ساز گرادیانی ممکن است روی فلات متوقف شود و
موفقیت گزارش کند؟ \textbf{پاسخ} قاعده‌ی توقفش بر اساس نزدیک به صفر بودن
گرادیان (یا اندازه‌ی گام) است. روی فلات گرادیان نزدیک صفر است با اینکه آن
نقطه کمینه نیست، پس معیار همگرایی برقرار می‌شود.

\textbf{6.} دو راه برای مثبت نگه داشتن یک پارامتر نرخ زوال در حین برازش
بگویید. \textbf{پاسخ} (1) استفاده از بهینه‌ساز کران‌دار (مثلاً
\lr{L-BFGS-B} یا \lr{Nelder\lr{--}Mead} با
\lr{\nolinkurl{bounds=[(1e-6, None)]}}). (2) \lr{reparameterization:}
\(\eta\) را روی کل اعداد حقیقی بهینه کنید و قرار دهید
\(\lambda = e^{\eta}\).

\textbf{7.} برازش‌های \lr{multi-start} مدل دوجزئی با
\((A_1,\allowbreak \lambda_1,\allowbreak A_2,\allowbreak \lambda_2) = (0.58,\allowbreak 0.36,\allowbreak 0.26,\allowbreak 0.00)\) و
\((0.26, 0.00, 0.58, 0.36)\) به یک \lr{SSE} کمینه‌ی یکسان می‌رسند. آیا این
مشکل بهینه‌ساز است؟ \textbf{پاسخ} نه. این جابه‌جایی برچسب \lr{(label
switching)} است: جابه‌جا کردن دو جزء دقیقاً همان منحنی را می‌دهد، پس مدل
بدون یک قید ترتیبی (مثلاً \(\lambda_1 > \lambda_2\)) \lr{identifiable}
نیست. بهینه‌ساز به‌درستی دو کمینه‌ی سراسری معادل را پیدا کرده است.

\begin{center}\rule{0.5\linewidth}{0.5pt}\end{center}

\subsection{تکلیف}

این هفته تکلیف نمره‌دار جدیدی منتشر نمی‌شود. \textbf{\lr{A1} (از کلامی به
فرمال) در ابتدای درس هفته‌ی 4 تحویل داده می‌شود.} \lr{lab} امروز همه‌ی
چیزهایی را که برای مرحله‌ی کالیبراسیون \lr{A1} لازم است به شما می‌دهد، اگر
بخواهید استخراج بسته‌ی خود را عددی بررسی کنید.

\textbf{تمرین اختیاری (بدون نمره، پیش از هفته‌ی 4 توصیه می‌شود):} تمرین‌های
3 و 4 \lr{lab} را کامل کنید. تمرین 3 (کمترین مربعات وزنی) پلی به بیشینه
درست‌نمایی است، و تمرین 4 پیش‌نمایشی از مسئله‌ی اصلی ارزیابی مدل است.
راه‌حل‌ها در ابتدای \lr{lab} هفته‌ی 4 به‌طور مختصر بحث می‌شوند.

\begin{center}\rule{0.5\linewidth}{0.5pt}\end{center}

\subsection{مطالعه‌ی بیشتر}

\begin{itemize}
\tightlist
\item
  \textbf{\lr{Farrell} \lr{\&} \lr{Lewandowsky} (2018)، فصل 3.} برآورد پایه‌ی
  پارامتر: توابع ناهمخوانی، سطوح خطا، الگوریتم سیمپلکس و مسائل عملی مثل
  کمینه‌های محلی و مقادیر شروع. \emph{(منبع اصلی.)}
\item
  \textbf{قدوسیان (1401)}، \emph{مبانی روش‌های مستقیم در بهینه‌سازی
  غیرخطی}: روش‌های جستجوی مستقیم برای بهینه‌سازی غیرخطی، از جمله روش‌های
  نوع سیمپلکس. برای دانشجویانی که نظریه‌ی پشت \lr{Nelder\lr{--}Mead} را
  می‌خواهند توصیه می‌شود {[}شماره‌ی فصل بررسی شود{]}.
\item
  \textbf{\lr{Verguts} (2022)}: فصل بهینه‌سازی و گرادیان نزولی، مفید برای
  نگاه گرادیانی که در هفته‌ی 15 (مدل‌های یادگیری) برمی‌گردد {[}شماره‌ی فصل
  بررسی شود{]}.
\end{itemize}

\emph{پیش‌نمایش هفته‌ی 4:} به‌جای پرسیدن «کدام پارامترها مربع خطا را کمینه
می‌کنند؟» می‌پرسیم «کدام پارامترها داده‌ی مشاهده‌شده را محتمل‌ترین می‌کنند؟»
این بیشینه درست‌نمایی است. داده‌های دوجمله‌ای را درست مدیریت می‌کند و پایه‌ی
\lr{AIC/BIC} (هفته‌ی 7) و استنباط بیزی (هفته‌های 9\lr{--}13) است.
